% !TEX root = ../main_ext_2.tex

\section*{Abstract}
\thispagestyle{empty}

Vector autoregressions (VARs) and local projections (LPs) estimate the same population impulse
responses, yet they are usually compared at a common lag order, which confounds the choice of
estimator with the choice of model complexity. We revisit this comparison in a Monte Carlo study
that holds effective complexity constant, following the complexity-adjusted benchmark of
\textcite{ludwig2026local}. Under correct specification, a fixed LP(4) already behaves like a
highly parameterized VAR of order $q \approx 15$, so the familiar point-estimation edge of
low-order VARs reflects parsimony rather than an intrinsic advantage of the iterated method. We
then introduce functional misspecification through a nonlinear lag-quadratic process. Because the
centered quadratic term is nearly orthogonal to the linear lag space, the target impulse response
stays almost unchanged, so the equal-lag VAR(4) keeps its root mean squared error advantage over
LP(4) and the finite-sample gap is driven by variance alone. The nonlinearity instead surfaces as
conditional heteroskedasticity in the residuals, which the homoskedastic VAR delta-method intervals
cannot accommodate. They undercover ever more severely as persistence and nonlinearity rise, while
the heteroskedasticity-robust intervals of the local projection stay near nominal. This is the first
of our designs in which the robustness of local projections becomes a tangible advantage, and it is
confined to inference. Population equivalence of the two estimands therefore does not carry over to
finite-sample inference, and the better method depends on whether the researcher prioritizes the
efficiency of the VAR point estimate or the interval validity of the direct projection.

\clearpage
